\documentclass[11pt]{article}
\usepackage[a4paper,margin=27mm]{geometry}
\usepackage{amsmath,amssymb,amsthm,mathtools,xcolor,booktabs,array}
\newcommand{\PROVED}{\textcolor{green!40!black}{\textbf{PROVED}}}
\newcommand{\OPEN}{\textcolor{red!70!black}{\textbf{OPEN}}}
\newcommand{\AUDIT}{\textcolor{orange!70!black}{\textbf{AUDIT}}}
\newtheorem{theorem}{Theorem}
\newtheorem{lemma}{Lemma}
\newtheorem{proposition}{Proposition}
\newtheorem{corollary}{Corollary}
\title{Excluding the almost-constant counterexample to URGT$_1$ (v125)}
\date{14 August 2026}

\begin{document}
\maketitle

\section{The first Gamma channel and the only dangerous band}

For $a_j=2j+\frac12$ put
\begin{equation}
 g_j(\tau)=\frac{2\tau^2}{a_j(a_j^2+\tau^2)},
 \qquad g_\Gamma(\tau)=\sum_{j\ge0}g_j(\tau).              \tag{1}
\end{equation}
The first channel, $a_0=\frac12$, is
\begin{equation}
 g_0(\tau)=\frac{4\tau^2}{\tau^2+\frac14}.                \tag{2}
\end{equation}

\begin{lemma}[Localization of every possible Gamma deficit --- \PROVED]
One has
\begin{equation}
 g_0(\tau)\ge1
 \quad\Longleftrightarrow\quad
 |\tau|\ge\frac1{\sqrt{12}}.                              \tag{3}
\end{equation}
Consequently, for every $f$ in the Gamma form domain,
\begin{equation}
 \mathcal G_\Gamma[f]-\|f\|^2
 \ge
 \int_{|\tau|<1/\sqrt{12}}
       (g_\Gamma(\tau)-1)|\widehat f(\tau)|^2
       \frac{d\tau}{2\pi}.                               \tag{4}
\end{equation}
In particular, every counterexample to URGT$_1$ must concentrate its
physical leakage in the fixed band
$|\tau|<1/\sqrt{12}$; no $N$-dependent enlargement of that band is
possible.
\end{lemma}

\begin{proof}
Equation (3) follows by multiplying (2) by
$\tau^2+\frac14>0$.  Outside the band in (3), the first summand alone is
at least one, and every remaining summand in (1) is nonnegative.  Dropping
the nonnegative exterior part gives (4).
\end{proof}

\section{Exact factorization of the depth-one rational polynomial}

Let
\begin{equation}
 A_N(\tau)=\sum_{n\le N}\frac{\Lambda(n)}{\sqrt n}\,n^{-i\tau}.
                                                                    \tag{5}
\end{equation}
For the depth-one coefficient in v124 define
\begin{equation}
 \Phi_{N,1}(\tau)=
 \sum_{a/b\in\mathscr R_N^{\rm leak}}
 \ell_{N,1}(a,b)(a/b)^{-i\tau}.                            \tag{6}
\end{equation}

\begin{proposition}[Full pair product minus the divisibility return ---
\PROVED]
One has the exact identity
\begin{equation}
 \boxed{
 \Phi_{N,1}(\tau)
 =\frac{|A_N(\tau)|^2}{\log N}-D_N(\tau),}                \tag{7}
\end{equation}
where
\begin{equation}
 D_N(\tau)=\frac1{\log N}
 \sum_{\substack{m,n\le N\\n\mid m}}
 \frac{\Lambda(m)\Lambda(n)}{\sqrt{mn}}
 \left(\frac mn\right)^{-i\tau}.                         \tag{8}
\end{equation}
For every fixed $R>0$,
\begin{equation}
 \sup_{|\tau|\le R}|D_N(\tau)|=O(\log N).                \tag{9}
\end{equation}
\end{proposition}

\begin{proof}
Expanding $|A_N(\tau)|^2$ gives all ordered pairs $(m,n)$ with weight
$\Lambda(m)\Lambda(n)/\sqrt{mn}$ and label $m/n$.  Aggregating equal
reduced fractions gives (2) of v124.  The label is an integer exactly when
$n\mid m$, which is (8).  This proves (7).

For (9), a nonzero summand in (8) has $n=p^r$ and $m=p^s$ for a common
prime $p$ and $1\le r\le s$.  Hence, uniformly in real $\tau$,
\[
 |D_N(\tau)|
 \le\frac1{\log N}\sum_{p^s\le N}(\log p)^2
       \sum_{1\le r\le s}p^{-(r+s)/2}.
\]
The stratum $r=s=1$ is $O((\log N)^2)$ by partial summation from the prime
number theorem.  The strata with $r+s\ge3$ are no larger after summing the
geometric tails.  Division by $\log N$ proves (9).
\end{proof}

\begin{corollary}[Compact-frequency continuum limit --- \PROVED]
For every fixed $R>0$,
\begin{equation}
 \boxed{
 \sup_{|\tau|\le R}
 \left|
 \frac{\log N}{N}\Phi_{N,1}(\tau)
 -\frac1{\tau^2+\frac14}
 \right|\longrightarrow0.}                               \tag{10}
\end{equation}
\end{corollary}

\begin{proof}
Uniform partial summation of the prime number theorem on compact
$\tau$-sets gives
\begin{equation}
 \frac{A_N(\tau)}{\sqrt N}
 =\frac{N^{-i\tau}}{\frac12-i\tau}+o_R(1).                 \tag{11}
\end{equation}
Therefore
$|A_N(\tau)|^2/N\to(\tau^2+\frac14)^{-1}$ uniformly on
$[-R,R]$.  Equations (7) and (9) show that the normalized divisibility
return is $O((\log N)^2/N)=o(1)$.  This proves (10).
\end{proof}

\section{The critical model pays exactly one unit}

Put $b=\frac12$ and
\begin{equation}
 W_*(\tau)=\frac1{(\tau^2+b^2)^2}.                         \tag{12}
\end{equation}
This is the squared continuum profile supplied by (10).

\begin{theorem}[Exact critical Gamma normalization --- \PROVED]
For every $a>0$,
\begin{equation}
 \frac{displaystyle
  \int_{\mathbb R}
  \frac{2\tau^2}{a(a^2+\tau^2)}W_*(\tau)\,d\tau}
 {\displaystyle\int_{\mathbb R}W_*(\tau)\,d\tau}
 =\frac{2b^2}{a(a+b)^2}.                                   \tag{13}
\end{equation}
Consequently
\begin{align}
 \frac{\int g_0W_*}{\int W_*}&=1,                         \tag{14}\\
 \frac{\int g_\Gamma W_*}{\int W_*}
 &=1+\eta_\Gamma,                                         \tag{15}\\
 \eta_\Gamma
 &:=\sum_{j\ge1}
 \frac{1/2}{(2j+\frac12)(2j+1)^2}
 >\frac1{45}.                                              \tag{16}
\end{align}
Thus the first Gamma channel pays the unit norm exactly on the unique
continuum profile suggested by a coherent depth-one leakage cloud.  The
second Gamma channel already supplies the strict surplus $1/45$; all
later channels add further positive energy.
\end{theorem}

\begin{proof}
First,
\begin{equation}
 \int_{\mathbb R}\frac{d\tau}{(\tau^2+b^2)^2}
 =\frac{\pi}{2b^3}.                                       \tag{17}
\end{equation}
Also
\[
 \int_{\mathbb R}
 \frac{\tau^2,d\tau}{(\tau^2+a^2)(\tau^2+b^2)}
 =\frac\pi{a+b}.
\]
Differentiating this identity with respect to $b$ gives
\begin{equation}
 \int_{\mathbb R}
 \frac{\tau^2,d\tau}
      {(\tau^2+a^2)(\tau^2+b^2)^2}
 =\frac\pi{2b(a+b)^2}.                                    \tag{18}
\end{equation}
Multiplication by $2/a$ and division by (17) proves (13).  Taking
$a=b=\frac12$ gives (14).  Substitution of
$a_j=2j+\frac12$ gives (15)--(16).  The first omitted term, $j=1$, is
\[
 \frac{1/2}{(5/2)3^2}=\frac1{45},
\]
and all subsequent terms are positive.
\end{proof}

\section{What the counterexample experiment now proves}

Let
\begin{equation}
 e_N(t)=\delta_N^{-1/2}\mathbf1_{(0,\delta_N)}(t).          \tag{19}
\end{equation}
On every fixed frequency interval,
$|\widehat e_N(\tau)|^2/\delta_N\to1$ uniformly.  Combining this with
(10) and Theorem 4 shows that the proposed almost-constant obstruction has
the critical Gamma ratio $1+\eta_\Gamma$, not a ratio below one.  A fully
rigorous passage over the whole real line uses (3): outside the fixed
deficit band the integrand of
$\mathcal G_\Gamma-\|\cdot\|^2$ is nonnegative, so only compact-frequency
convergence is needed for the lower bound.

\begin{corollary}[No depth-one almost-constant counterexample ---
\PROVED]
For the profiles (19),
\begin{equation}
 \liminf_{N\to\infty}
 \frac{
 \mathcal G_\Gamma[L_{N,1}e_N]
 -\|L_{N,1}e_N\|^2}
 {\delta_NN^2/(\log N)^2}
 \ge
 \eta_\Gamma
 \int_{\mathbb R}W_*(\tau)\frac{d\tau}{2\pi}>0.          \tag{20}
\end{equation}
In particular, the coherent depth-one leakage with a constant cell
profile satisfies the Gamma domination strictly for all sufficiently
large $N$, even before adding the Tate terminal.
\end{corollary}

\begin{proof}
Fix $R>1/\sqrt{12}$.  On $[-R,R]$, use (10) and the uniform convergence
$|\widehat e_N|^2/\delta_N\to1$.  Outside the deficit band, the integrand
$g_\Gamma-1$ is nonnegative by (3), so it may be discarded when taking a
lower bound.  Let first $N\to\infty$ and then $R\to\infty$.
Equations (15)--(16) give (20).
\end{proof}

\paragraph{Scope of the exclusion.}
Corollary 5 eliminates the most dangerous proposed counterexample and
explains the unit normalization.  It is not yet the full URGT$_1$ theorem,
because an arbitrary cell profile may have zero mean, and all Witt depths
must be handled simultaneously.  A remaining counterexample would have
to satisfy all three conditions:
\begin{enumerate}
 \item it is asymptotically orthogonal to the constant cell profile;
 \item its leakage remains concentrated in $|\tau|<1/\sqrt{12}$ despite
       support of the source in a cell of length $\delta_N\asymp N^{-1}$;
 \item the higher-depth polynomials cancel the strict surplus (16), even
       though depths are orthogonal and every Gamma summand is positive.
\end{enumerate}
The next theorem to prove is therefore the uniform mean-zero stability
estimate
\begin{equation}
 \boxed{
 \sum_{k\ge1}\int_{|\tau|<1/\sqrt{12}}
 (1-g_\Gamma(\tau))
 |\widehat e_0(\tau)|^2|\Phi_{N,k}(\tau)|^2\frac{d\tau}{2\pi}
 \le
 \sum_{k\ge1}\int_{|\tau|\ge1/\sqrt{12}}
 (g_\Gamma(\tau)-1)
 |\widehat e_0(\tau)|^2|\Phi_{N,k}(\tau)|^2\frac{d\tau}{2\pi}
 +\mathcal T_N(e_0),}                                      \tag{21}
\end{equation}
where $\int_0^{\delta_N}e_0=0$.  Unlike v124 (20), (21) has removed the
only saturating scalar mode.

\begin{center}
\begin{tabular}{>{\raggedright\arraybackslash}p{.57\linewidth}
                >{\centering\arraybackslash}p{.14\linewidth}
                >{\raggedright\arraybackslash}p{.19\linewidth}}
\toprule
Statement & Status & Location \\
\midrule
Fixed dangerous frequency band & \PROVED & (2)--(4) \\
Exact depth-one product/divisibility split & \PROVED & (7)--(9) \\
Continuum leakage profile & \PROVED & (10)--(11) \\
First Gamma channel pays exactly one & \PROVED & (13)--(14) \\
Strict higher-Gamma margin $>1/45$ & \PROVED & (15)--(16) \\
No almost-constant depth-one counterexample & \PROVED & (19)--(20) \\
Uniform mean-zero stability, all depths & \OPEN & (21) \\
URGT$_1$, $G$, row (d) & \OPEN & after (21) and finite audit \\
\bottomrule
\end{tabular}
\end{center}

\end{document}
